% ============================== CI-8 ==============================
\reading{CI-8}{Digital Resilience and Financial Stability: The Quest for Policy Tools in the Financial Sector}{DEFER}{0}

\begin{cikeybox}[title={The three objectives, in the spine's own words}]
\begin{enumerate}[label=\textbf{\alph*.}]
\item Describe characteristics of cyber risks and information/communication technology (ICT)
risks faced by financial institutions.
\item Assess the interactions between cyber and ICT risks and financial risks and explain how
cyber and ICT risk events at financial institutions can lead to systemic financial risk.
\item Describe potential macroprudential tools and policy measures that can be used to address
cyber risks and ICT risks and explain challenges to the adoption of each one.
\end{enumerate}
\greyline{Source: Jose Ramon Martinez, ``Digital Resilience and Financial Stability, The Quest
for Policy Tools in The Financial Sector'', Banco de Espa\~na (April 2023). The class notes
label this Reading 107; it is CI-8 on the 2026 spine.}
\end{cikeybox}

\begin{cisrcbox}
{\sffamily\bfseries\color{FRMNavy}A structural defect in the source notes, corrected here.}
The notes jump straight from LO 107.a to LO 107.c. Objective (b) has no heading anywhere in
them. The material is not missing, though: the block headed ``Interactions Between Cyber Risks
and Financial Stability'', with its five numbered points, sits inside the 107.a cell with no
label, and it is a point-for-point match to objective (b).\par\vspace{4pt}
So this is a \textbf{reattribution, not a gap-fill}. The body text below is the notes' own
words, moved to the objective it answers and given the spine's heading. No gold badge appears,
because nothing here was written from an outside source. The match was confirmed against
SchweserNotes Book 5, LO 107.b.
\end{cisrcbox}

% ------------------------------------------------------------------
\lo{CI-8 a}{Describe characteristics of cyber risks and information/communication technology (ICT) risks faced by financial institutions.}

\begin{center}
\renewcommand{\arraystretch}{1.25}
\begin{tabularx}{\textwidth}{@{}p{2.1cm} >{\raggedright\arraybackslash}X
                              >{\raggedright\arraybackslash}X@{}}
\toprule
\rowcolor{FRMSoft}
&{\sffamily\bfseries\color{FRMRed}Cyber risks} &
{\sffamily\bfseries\color{FRMNavy}ICT risks}\\
\midrule
\term{Definition} &
Loss of confidentiality, integrity, and availability (CIA) of data and data systems
\emph{resulting from cyberattacks} &
Disruptions in information and communication systems that \emph{may not stem from
cyberattacks}\\
\addlinespace
\term{Examples} &
Breaches of client information, unauthorised access, funds theft, DoS attacks &
System breakdowns, communication failures, service outages\\
\addlinespace
\term{Importance} &
Ranked as significant as market risk by some financial institutions &
Essential for ensuring uninterrupted transaction services\\
\addlinespace
\term{Challenges} &
Difficulty in estimating losses due to lack of reliable data, which has hindered the
development of cyber insurance &
Dependency on technology, vulnerability to system failures\\
\bottomrule
\end{tabularx}
\end{center}
\figcap{Table CI-8.1. The two are told apart by \emph{cause}, not by consequence. A service
outage looks identical to the customer whether it came from an attacker or from a failed
upgrade, but only the first is cyber risk. That is also why the challenge rows differ: cyber
risk suffers from missing loss data because attacks are under-reported, while ICT risk suffers
from plain technological dependency.}

\subsection*{Cloud computing}

Cloud computing benefits financial institutions, but comes with major risks:

\begin{enumerate}[label=\textbf{\alph*)}]
\item Shared cyber risks with cloud providers.
\item Widespread adoption increases systemic cyberattack risk.
\item Multiclouds reduce single-provider risk but introduce complexity.
\item Frequent outages raise concerns about potential disruptions.
\end{enumerate}

% ------------------------------------------------------------------
\lo{CI-8 b}{Assess the interactions between cyber and ICT risks and financial risks and explain how cyber and ICT risk events at financial institutions can lead to systemic financial risk.}

Cyber and ICT risks at financial institutions can trigger systemic financial risks, impacting
\term{liquidity}, \term{leverage}, and \term{customer trust}.

\begin{enumerate}
\item \term{Integration and contagion effects.} The highly integrated nature of cyber and ICT
connections within financial institutions creates the potential for significant contagion
effects.
\item \term{Challenges in predicting probability.} There is difficulty in predicting key events
due to nonlinear contagion effects. The response is to use \term{agent-based models (ABMs)} for
simulation. ABMs assess impacts on confidentiality, integrity, and availability, which
influence liquidity and asset values.
\item \term{Tracing systemic risks.} Cyberattacks might block direct actions like bank
withdrawals but can still cause financial contagion indirectly. ABMs can be used here as well,
simulating individual behaviours to identify and understand hidden and systemic risks.
\item \term{Reliance on digital infrastructure.} Financial institutions depend heavily on
digital infrastructure, including cloud computing. Quick recovery from cyber incidents is
essential to maintain financial stability and continuous operation.
\item \term{Impact on customer confidence.} Cyber risks can diminish customer confidence,
potentially leading to bank runs and systemic financial instability.
\end{enumerate}

\begin{cidefbox}[title={CIA triad}]
Stands for \term{Confidentiality}, \term{Integrity}, and \term{Availability}, the three key
principles of cybersecurity. The impact of cyber incidents on these can vary: \term{integrity
and availability are often more critically affected than confidentiality}.
\end{cidefbox}

\begin{cifigblock}
{\centering
\begin{tikzpicture}[
  font=\sffamily\scriptsize,
  ev/.style={draw=FRMRule, fill=FRMSoft, text width=2.5cm, align=center,
             minimum height=1.0cm, inner sep=3pt, rounded corners=1pt},
  cia/.style={draw=FRMRule, text width=2.5cm, align=center,
              minimum height=0.75cm, inner sep=3pt, rounded corners=1pt},
  fin/.style={draw=FRMRule, fill=PaleGold, text width=2.5cm, align=center,
              minimum height=0.75cm, inner sep=3pt, rounded corners=1pt},
  sys/.style={draw=FRMRule, fill=PaleRed, text width=2.6cm, align=center,
              minimum height=1.0cm, inner sep=3pt, rounded corners=1pt},
  ar/.style={-{Stealth[length=4pt]}, FRMRed, line width=0.8pt}
]
\node[ev] (e) at (0,0) {cyber or ICT event};

\node[cia, fill=PaleRed!55] (c1) at (3.4,1.15) {\textbf{Integrity}};
\node[cia, fill=PaleRed!55] (c2) at (3.4,0.15) {\textbf{Availability}};
\node[cia, fill=PaleNavy] (c3) at (3.4,-0.95) {\textbf{Confidentiality}};
\node[above=2.5mm of c1, font=\sffamily\bfseries\scriptsize\color{FRMNavy}] {CIA triad};
\node[below=2.5mm of c3, font=\sffamily\scriptsize\itshape\color{FRMGrey}, text width=3.6cm,
      align=center]
  {shaded red: more critically affected than confidentiality};

\node[fin] (f1) at (8.05,1.15) {liquidity};
\node[fin] (f2) at (8.05,0.15) {leverage};
\node[fin] (f3) at (8.05,-0.95) {customer trust};
\node[above=2mm of f1, font=\sffamily\bfseries\scriptsize\color{FRMNavy}]
  {financial risks};

\node[sys] (s) at (11.3,0.15) {systemic financial risk};

\draw[ar] (e.east) -- (c1.west);
\draw[ar] (e.east) -- (c2.west);
\draw[ar] (e.east) -- (c3.west);
\draw[ar] (c1.east) -- ++(0.55,0) -- (f1.west);
\draw[ar] (c2.east) -- ++(0.55,0) -- (f2.west);
\draw[ar] (c3.east) -- ++(0.55,0) -- (f3.west);
\draw[ar] (f1.east) -- (s.west);
\draw[ar] (f2.east) -- (s.west);
\draw[ar] (f3.east) -- (s.west);

\node[font=\sffamily\scriptsize\itshape\color{FRMGrey}, text width=10.0cm, align=center]
  at (5.7,-3.05)
  {contagion effects are nonlinear, so the probability of the next step cannot be predicted
   directly: agent-based models simulate it instead};
\end{tikzpicture}\par}
\figcap{Figure CI-8.1. The objective is about the join between two worlds, and the join runs
through the CIA triad. A technology event lands on confidentiality, integrity or availability,
those land on liquidity, leverage and customer trust, and only then does it become systemic.
The ordering inside the triad is the examinable detail and it is counter-intuitive: it is
integrity and availability, not confidentiality, that are more critically affected. The arrows
are drawn one-to-one for legibility; the notes do not pair specific CIA components with
specific financial risks.}
\end{cifigblock}

% ------------------------------------------------------------------
\lo{CI-8 c}{Describe potential macroprudential tools and policy measures that can be used to address cyber risks and ICT risks and explain challenges to the adoption of each one.}

\begin{cidefbox}[title={Digital resilience}]
Protecting financial systems from cyber and ICT risks requires a two-part approach:
\begin{enumerate}
\item \term{Micro-oriented tools}: operational guidance for individual entities.
\item \term{Framework policies}: fundamental rules for network and system operations.
\end{enumerate}
\end{cidefbox}

\subsection*{Policy measures and micro-oriented tools}

\begin{enumerate}
\item \term{International Telecommunications Union (ITU)}: implements cyber resilience
measures; emphasises certification and standardisation for effective cybersecurity.
\item \term{European Union (EU)}:
  \begin{itemize}
  \item established policies for incident reporting and information sharing;
  \item the \term{Digital Operational Resilience Act (DORA)} focuses on ICT third-party risks,
  including cloud computing, and mandates penetration stress tests to improve ongoing security
  awareness.
  \end{itemize}
\end{enumerate}

\subsection*{Macroprudential tools}

\begin{cifigblock}
{\centering
\begin{tikzpicture}[
  font=\sffamily\scriptsize,
  tool/.style={draw=FRMRule, fill=PaleNavy, text width=4.4cm, align=left, inner sep=5pt,
               rounded corners=1pt},
  chal/.style={draw=FRMRule, fill=PaleRed, text width=4.4cm, align=left, inner sep=5pt,
               rounded corners=1pt, font=\sffamily\scriptsize\itshape}
]
\node[tool] (t1) at (0,0)
  {\textbf{1. Circuit breakers}
   \begin{itemize}[leftmargin=10pt, topsep=2pt, itemsep=1pt]
   \item temporary pauses on normal rules and regulations, a temporary stoppage to some
   services
   \item in a major cyber incident affecting banks, they minimise the negative impacts
   \end{itemize}};
\node[tool, below=3mm of t1] (t2)
  {\textbf{2. Cooperative arrangements}
   \begin{itemize}[leftmargin=10pt, topsep=2pt, itemsep=1pt]
   \item cooperation among financial institutions is essential for robust defences
   \item information and knowledge-sharing resources support that cooperation
   \item \textbf{ICT buffers} sit here: co-owner banks hit by an attack can let customers
   transact on their accounts at unaffected banks
   \end{itemize}};
\node[tool, below=3mm of t2] (t3)
  {\textbf{3. Structural measures}
   \begin{itemize}[leftmargin=10pt, topsep=2pt, itemsep=1pt]
   \item regulating systemically important technological institutions (SITIs)
   \item ``forced diversification'': using multiple cloud providers
   \item penalties for concentrated software use
   \end{itemize}};

\node[chal, right=8mm of t1] (c1)
  {Challenge: defining the severity of a cyberattack that would warrant activating the
   circuit breaker};
\draw[-{Stealth[length=4pt]}, FRMRed, line width=0.7pt] (t1.east) -- (c1.west);
\end{tikzpicture}\par}
\figcap{Figure CI-8.2. Three tools, and note where the ICT buffer actually sits: it is part of
the \emph{cooperative} measures, not a structural one, even though forced diversification in
the third box solves a neighbouring problem. The challenge the notes state explicitly belongs
to the circuit breaker, and it is a threshold problem rather than a design problem: the tool
works, but nobody has defined how bad an attack has to be before it fires.}
\end{cifigblock}

\begin{cikeybox}[title={Where the volume ends}]
ICT buffers are part of collaborative measures aimed at combating systemic risks and ensuring
rapid recovery from major cyberattacks. Co-owner banks affected by cyberattacks can allow
customers to transact on their accounts at unaffected banks, reducing the overall risk.
\end{cikeybox}
